% Generated supplemental LaTeX fragment with expanded cheatsheet content
\section{Bővített Cheatsheet tartalom}

A következő rész a \texttt{CHEATSHEET.md}-ben található fontos parancsok bővített, magyarázó változata.

\subsection{Fájlok törlése tulajdonos szerint}
\begin{lstlisting}[language=bash,caption={Fájlok törlése user szerint}]
# List for verification:
find /path/to/dir -maxdepth 1 -type f -user operator -print

# Interactive delete (recommended):
find /path/to/dir -maxdepth 1 -type f -user operator -exec rm -i {} \;

# Automatic, dangerous delete:
find /path/to/dir -maxdepth 1 -type f -user operator -exec rm -f {} \;
\end{lstlisting}

Példa output (lista):
\begin{verbatim}
/path/to/dir/file1.txt
/path/to/dir/log-2026-05-01.log
/path/to/dir/.secret
\end{verbatim}

\subsection{RSA kulcsok generálása és PKCS\#8 konverzió}
A következő példák bemutatják, hogyan generáljunk 512 bites RSA kulcsokat és hogyan alakítsuk át őket PKCS\#8 formátumba.

\begin{lstlisting}[language=bash,caption={RSA 512 példa}]
mkdir -p dev/example
openssl genrsa -out dev/example/rsa-512-01.pem 512
openssl pkcs8 -topk8 -inform PEM -outform PEM -nocrypt -in dev/example/rsa-512-01.pem -out dev/example/rsa-512-01-pkcs8.pem
\end{lstlisting}

További hasonló fájlok létrehozása automatizáltan:
\begin{lstlisting}[language=bash]
for i in {2..11}; do
  openssl genrsa -out dev/example/rsa-512-0${i}.pem 512 && \
  openssl pkcs8 -topk8 -inform PEM -outform PEM -nocrypt -in dev/example/rsa-512-0${i}.pem -out dev/example/rsa-512-0${i}-pkcs8.pem
done
\end{lstlisting}

Példa (kivonat a PKCS\#8 fájlból):
\begin{verbatim}
MIIEvgIBADANBgkqhkiG9w0BAQEFAASCBK...
\end{verbatim}

Megjegyzés: 512 bit demonstrációs célokra; ne használjuk erős titkosításhoz.

\subsection{Cron kalkulátor link és példa}
Használjunk megbízható cron kalkulátort a feladatok ütemezésének ellenőrzéséhez: \url{https://crontab.guru/}

Példa cron bejegyzés és magyarázat:
\begin{lstlisting}
# Runs every day at 02:00:
0 2 * * * /usr/local/bin/backup.sh
# Cron calculator: https://crontab.guru/0%202%20*%20*%20*
\end{lstlisting}

\subsection{Output példák a parancsokra}
\begin{itemize}
  \item \texttt{lsblk -f} részlet:
\begin{verbatim}
NAME   FSTYPE LABEL UUID                                 MOUNTPOINT
sda                                                          
|--sda1 ext4         3a5f-...                           /
`--sda2 swap         8f3b-...                           [SWAP]
\end{verbatim}

  \item \texttt{openssl rsa -in key.pem -pubout} kimenet (részlet):
\begin{verbatim}
-----BEGIN PUBLIC KEY-----
MIIBIjANBgkqhkiG9w0BAQEFAAOCAQ8AMIIBCgKCAQEAv...
-----END PUBLIC KEY-----
\end{verbatim}

\end{itemize}

\subsection{Hivatkozások}
Bővebb olvasmányok és referencia oldalak:
\begin{itemize}
  \item Cron kifejezések: \url{https://crontab.guru/}
  \item OpenSSL dokumentáció: \url{https://www.openssl.org/docs/man1.1.1/man1/}
  \item find(1) man: \url{https://linux.die.net/man/1/find}
\end{itemize}
